\section{Conclusão}\label{sec:conclusao}

Este trabalho apresentou o \texttt{mcp-kicad-vscode}, um servidor Model Context
Protocol para automação do fluxo de projeto de placas de circuito impresso no
KiCad. A solução expõe 233 ferramentas organizadas em 24 módulos, 173 delas
descobríveis em 16 categorias, além de 23 recursos dinâmicos de estado do
projeto, integrações com JLCPCB, LCSC e DigiKey e suporte ao roteador
automático Freerouting. A validação combinou métricas de implementação,
execução das suítes de teste --- 2.412 de 2.451 casos aprovados localmente, com
a suíte TypeScript integralmente aprovada --- e um estudo de caso com uma placa
seguidora de linhas de quatro camadas (97 \emph{footprints}, 83 redes), cuja
verificação de regras reportou 11 violações residuais e zero itens não
conectados.

Os resultados respondem à questão de pesquisa na medida em que demonstram ser
viável, com uma arquitetura em duas camadas e disciplina de testes, expor um
fluxo completo de EDA a assistentes de IA de forma padronizada e verificável.
As limitações --- integração IPC não exercitada localmente, cobertura de
descoberta parcial e avaliação restrita a um único projeto --- delimitam o
alcance das conclusões e orientam a agenda de continuidade, que inclui a
validação em tempo real com a interface gráfica, a ampliação da cobertura de
testes e a realização de estudos com usuários.
